% 数列の超体系的解説：付録改訂版
% 漸化式から、連続・反復・近似・母関数・行列・多項式・漸近解析へ進む。

\chapter{付録}

ここまで本編では、漸化式を見たときにどのような変形をすれば一般項にたどり着けるかを考えてきた。
付録では、少しだけ歩く方向を変えてみよう。

漸化式は、項を順番に計算するための規則である。
しかしそれだけではない。
「次の値を決める」という考え方は、微分方程式、数値計算、ニュートン法、母関数、行列、三角関数、多項式、そして極限の理論へとつながっている。

ここで目指すのは、それぞれの理論を詳しく学び尽くすことではない。
漸化式を出発点にして、数学の別の場所へどうやって歩いていけるのかを眺めることである。
式を暗記するのではなく、「なぜその式を考えたくなったのか」を大切にして進もう。

なお、この付録では話の流れに合わせて添字を $0$ から始めることがある。
その場合は、本文の添字の付け方とは別に、各節で明示する。

\section{まず、変化を記録する}

\subsection{漸化式は変化の記録である}

数列
\[
  a_0,a_1,a_2,\ldots
\]
を見たとき、各項の値だけを並べることもできる。
しかし、数列がどのように変化しているかを知りたいなら、隣り合う項の差を調べる方が自然である。

\[
  \Delta a_n=a_{n+1}-a_n
\]

この $\Delta a_n$ を数列の\textbf{差分}という。
例えば、
\[
  a_n=n^2
\]
ならば、
\[
  \Delta a_n=(n+1)^2-n^2=2n+1
\]
である。
さらにもう一度差分を取ると、
\[
  \Delta^2a_n
  =\Delta(2n+1)=2
\]
となる。

ここで興味深いことが起こる。
二次式は、二回差分を取ると定数になるのである。
これは、関数において二回微分すると二次式が定数になることとよく似ている。

つまり、差分は「離散的な世界で変化の速さを見る道具」になっている。
この見方から、次の二つの問いが生まれる。

\begin{itemize}
  \item 差分と微分は、どの程度似ているのだろうか。
  \item 数列を、値を順番に生み出す仕組みとして見たら何が見えるだろうか。
\end{itemize}

この二つの問いを順番に追っていく。

\subsection{漸化式を「繰り返し」として見る}

一項前から次の項を決める漸化式が、
\[
  a_{n+1}=f(a_n)
\]
という形になっているとする。
これは、関数 $f$ を一回適用するたびに次の値へ進む規則である。

\[
  a_0
  \longrightarrow f(a_0)
  \longrightarrow f(f(a_0))
  \longrightarrow \cdots
\]

このように同じ関数を何度も適用することを\textbf{反復}という。
したがって、
\[
  a_n=f^{\,n}(a_0)
\]
と書ける。
ここで $f^{\,n}$ は $f$ を $n$ 回適用するという意味であり、$f(x)$ の $n$ 乗ではない。

この見方をすると、数列の問題は「値を計算する問題」から「関数を繰り返したとき、値がどこへ進むかを調べる問題」へ変わる。
例えば、ある値 $L$ が
\[
  f(L)=L
\]
を満たすなら、そこへ到達した後は動かない。
このような $L$ を\textbf{不動点}という。

不動点は、後で極限を考えるときに重要になる。
ただし、不動点があるからといって、どんな初期値からでもそこへ近づくとは限らない。
この「固定される場所」と「実際の動き」の違いも、付録の後半で確認する。

\subsection{ここから先の道筋}

差分を調べると、微分との類似が見えてくる。
反復として見ると、関数の不動点や近似の問題が見えてくる。
数列を一つの式にまとめれば母関数になり、複数の数列を一つのベクトルにまとめれば行列になる。

付録では、次の順に進む。

\[
\begin{array}{c}
\text{差分}\\[-2pt]
\downarrow\\[-2pt]
\text{微分方程式との対応}\\[-2pt]
\downarrow\\[-2pt]
\text{反復と近似}\quad\text{（ニュートン法）}\\[-2pt]
\downarrow\\[-2pt]
\text{数列をまとめる}\quad\text{（母関数・行列）}\\[-2pt]
\downarrow\\[-2pt]
\text{二項間漸化式と三角関数}\quad\text{（チェビシェフ多項式）}\\[-2pt]
\downarrow\\[-2pt]
\text{大きな }n\text{ での振る舞い}\quad\text{（漸近解析）}
\end{array}
\]

順番に別の理論を覚えるのではなく、同じ漸化式を別の角度から見直していると思ってほしい。

\section{離散から連続へ}

\subsection{差分と微分}

関数 $x(t)$ の微分は、
\[
  \frac{dx}{dt}
  =\lim_{h\to0}\frac{x(t+h)-x(t)}{h}
\]
で定義される。
これは、時間をほんの少しだけ進めたときの変化を、時間の幅で割ったものである。

一方、数列では、時刻を少しずつではなく一段ずつ進める。
したがって、
\[
  a_{n+1}-a_n
\]
が、離散的な変化を記録する。

微分と差分は同じものではない。
微分は極限を取った連続的な量であり、差分は実際の一段分の変化である。
それでも、
\[
  \frac{x(t+h)-x(t)}{h}
  \quad\longleftrightarrow\quad
  \frac{a_{n+1}-a_n}{h}
\]
という対応を考えると、両者を結ぶ道が見えてくる。

ここで $h$ は一段の時間幅である。
後で $h$ を小さくしていくと、差分から微分へ近づいていく。

\subsection{微分方程式から漸化式を作る}

物体の位置を表す関数 $x(t)$ が、
\[
  \frac{dx}{dt}=F(x)
\]
を満たしているとする。
この式は、「その瞬間の変化の速さが、現在の位置 $x$ によって決まる」という意味である。

この方程式をそのまま解くのが難しいとき、時刻を一定の幅 $h$ で区切る。
\[
  t_n=nh,\qquad a_n=x(t_n)
\]
とおくと、微分を前進差分で近似して、
\[
  \frac{a_{n+1}-a_n}{h}\approx F(a_n)
\]
となる。
従って、
\[
  a_{n+1}=a_n+hF(a_n)
\]
という漸化式を得る。

この方法は\textbf{オイラー法}と呼ばれる。
名前を覚えることが目的ではない。
「連続的な変化を、一段ずつの更新に置き換える」という発想が大切である。

\subsection{例：冷めていく物体を一段ずつ追う}

温度や濃度が、現在の値との差に比例して変化する場面を考えよう。
最も簡単なモデルとして、
\[
  \frac{dx}{dt}=1-x,\qquad x(0)=0
\]
を考える。
$x$ は時間が経つと $1$ に近づく。

時間幅を $h$ としてオイラー法を使うと、
\[
  a_{n+1}=a_n+h(1-a_n)=(1-h)a_n+h
\]
となる。
例えば $h=1/2$ なら、
\[
  a_0=0,\quad a_1=\frac12,\quad
  a_2=\frac34,\quad a_3=\frac78
\]
である。

この数列は、連続的な運動を完全に再現しているわけではない。
ただし、幅 $h$ を細かくすれば、より細かく変化を追える。
「解を求められないから終わり」ではなく、「一歩ずつ近づく数列を作る」という考え方がここで現れる。

\subsection{漸化式から連続的なモデルを想像する}

逆に、
\[
  a_{n+1}=a_n+hF(a_n)
\]
という漸化式を見たら、
\[
  \frac{a_{n+1}-a_n}{h}=F(a_n)
\]
であるから、
\[
  \frac{dx}{dt}=F(x)
\]
という微分方程式を想像できる。

ここで大切なのは、これは同値変形ではないということである。
漸化式は一段ごとの正確な規則であり、微分方程式はその動きを連続的に眺めたモデルである。
特に、元の漸化式に $h$ が書かれていない場合は、時間幅を $1$ としたと解釈していることになる。

連続化は、増加・減少・安定性を大づかみに考えるための道具である。
離散的な振動などを見落とすこともあるので、最後は元の漸化式へ戻って確認しなければならない。

\subsection{不動点だけでは収束しない}

例えば、
\[
  a_{n+1}=2-a_n
\]
を考える。
不動点 $L$ は、
\[
  L=2-L
\]
より $L=1$ である。
しかし $a_0=0$ とすると、
\[
  0,2,0,2,0,2,\ldots
\]
となり、$1$ へ近づかない。

「極限があるなら不動点でなければならない」ということと、「不動点だから極限になる」ということは別である。
この区別は、漸化式を連続的なモデルで考えるときにも、常に意識しておきたい。

\section{解そのものではなく、解に近づく数列を作る}

\subsection{方程式を漸化式に変える}

方程式
\[
  f(x)=0
\]
の解を直接求めたい。
しかし、解の公式が使えない方程式もある。
そのとき、「解を一発で求める」のではなく、「解に近づく数を順番に作る」ことを考える。

つまり、
\[
  x_0,x_1,x_2,\ldots
\]
を作り、
\[
  x_n\longrightarrow \alpha,\qquad f(\alpha)=0
\]
となることを目指す。
この時点で、方程式は漸化式の問題へ姿を変えている。

\subsection{接線を使うという発想}

いま $x=a$ の近くで $f(x)$ の値を知りたいとする。
曲線そのものは複雑でも、そこを通る接線なら扱いやすい。
テイラー展開の最初の二項だけを使うと、
\[
  f(x)\approx f(a)+f'(a)(x-a)
\]
となる。

接線が $x$ 軸と交わる点を求めるために、右辺を $0$ とおく。
\[
  0=f(a)+f'(a)(x-a)
\]
を $x$ について解くと、
\[
  x=a-\frac{f(a)}{f'(a)}
\]
を得る。

ここで $a=x_n$ と置けば、
\[
  \boxed{
  x_{n+1}=x_n-\frac{f(x_n)}{f'(x_n)}
  }
\]
となる。
これが\textbf{ニュートン法}である。

公式を覚えるより、
\[
  \text{曲線を接線で近似する}
  \longrightarrow
  \text{接線の交点を次の値にする}
\]
という流れを理解しておく方が大切である。

\subsection{例：平方根を漸化式で求める}

方程式
\[
  x^2-2=0
\]
の正の解は $\sqrt2$ である。
これをニュートン法で近似する。

\[
  f(x)=x^2-2,\qquad f'(x)=2x
\]
だから、
\[
  x_{n+1}
  =x_n-\frac{x_n^2-2}{2x_n}
  =\frac{x_n^2+2}{2x_n}
\]
となる。
初期値を $x_0=1$ とすると、
\[
  x_1=\frac32,\qquad
  x_2=\frac{17}{12},\qquad
  x_3=\frac{577}{408}
\]
であり、
\[
  1.414\ldots
\]
へ急速に近づいていく。

この場合は、$x_n>0$ のとき、誤差について
\[
  x_{n+1}-\sqrt2
  =\frac{(x_n-\sqrt2)^2}{2x_n}
\]
が成り立つ。
次の誤差が、前の誤差の二乗を含む形になっているため、誤差が急速に小さくなるのである。

\subsection{いつでもうまくいくわけではない}

ニュートン法は強力だが、どんな初期値からでも必ず成功するわけではない。
$f'(x_n)=0$ になると公式が使えないし、接線が遠くへ飛ぶこともある。

例えば、
\[
  f(x)=x^3-2x+2
\]
に $x_0=0$ からニュートン法を適用すると、
\[
  x_1=1,\qquad x_2=0
\]
となり、$0$ と $1$ を繰り返してしまう。

ここから分かるのは、漸化式を作っただけでは十分でないということだ。
その漸化式が、どの初期値から、どの範囲で、どのように動くのかを調べる必要がある。
これは本編で扱った漸化式にも共通する態度である。

\section{複雑な式を、よく知った形へ変える}

\subsection{変数を変える理由}

漸化式を解くとき、元の数列 $a_n$ をそのまま追わず、
\[
  b_n=g(a_n)
\]
という新しい数列を考えることがあった。
これは、難しい式を無理に計算するのではなく、見方を変えて簡単な規則を見つけるためである。

微分方程式でも、まったく同じ発想が現れる。
ここではその代表例として、リカッチ型方程式を眺める。
この節は本編の直接の解法に必要な知識ではないが、「変数変換」という考え方がどこまで広がるかを示す例になっている。

\subsection{リカッチ型方程式}

\[
  \frac{dx}{dt}=a(t)x^2+b(t)x+c(t)
\]
の形の微分方程式を\textbf{リカッチ型方程式}という。
$x^2$ が現れるので、$x$ について非線形である。

もし特別な解 $x_0(t)$ が一つ分かっているとする。
すると、未知の解 $x(t)$ と既知の解 $x_0(t)$ の差を利用したくなる。
そこで、
\[
  x=x_0+\frac1u
\]
と置く。
これは「解そのもの」ではなく、「既知の解からどれだけずれているか」を $u$ で記録する変換である。

微分すると、
\[
  \frac{dx}{dt}=\frac{dx_0}{dt}-\frac1{u^2}\frac{du}{dt}
\]
となる。
これを元の方程式へ代入し、$x_0$ が元の方程式を満たすことを使うと、
\[
  \frac{du}{dt}+\{2a(t)x_0(t)+b(t)\}u=-a(t)
\]
を得る。

右辺に $u$ の二乗がない。
つまり、非線形な方程式を、$u$ についての線形方程式へ変えることができた。
「分かっているものを基準にして、未知の部分だけを別の変数で表す」という発想が、ここで働いている。

\subsection{小さな例で変換を確認する}

\[
  \frac{dx}{dt}=x^2-1
\]
を考える。
$x_0(t)=1$ は定数解である。
\[
  x=1+\frac1u
\]
とおくと、
\[
  -\frac1{u^2}\frac{du}{dt}
  =\left(1+\frac1u\right)^2-1
  =\frac2u+\frac1{u^2}
\]
である。
両辺に $u^2$ を掛けると、
\[
  \frac{du}{dt}+2u=-1
\]
となる。

初めは $x^2$ を含む非線形方程式だったものが、最後には高校数学で扱う線形の方程式になった。
数列でも微分方程式でも、式を変形する目的は「計算を増やすこと」ではなく、「既知の形へ移ること」である。

\section{数列を一つの式にまとめる}

\subsection{母関数が欲しくなる理由}

数列は無限個の数を持っている。
そのため、項を一つずつ計算していると、全体の構造が見えにくいことがある。
そこで、各項を係数として一つの式に詰め込む。

数列 $a_0,a_1,a_2,\ldots$ に対して、
\[
  A(x)=a_0+a_1x+a_2x^2+\cdots
\]
を考える。
これを数列の\textbf{母関数}という。

ここでは、まず収束する関数としてではなく、係数を比較するための\textbf{形式的冪級数}として扱う。
つまり、$x$ に具体的な数を代入することよりも、$x^n$ の係数を読むことが目的である。

数列を母関数へ移すことで、
\[
  (a_0,a_1,a_2,\ldots)
  \longrightarrow A(x)
\]
という視点の変更が起こる。

\subsection{フィボナッチ型漸化式を母関数へ移す}

\[
  a_{n+2}=a_{n+1}+a_n
  \qquad(n\geq0)
\]
を考える。
母関数
\[
  A(x)=a_0+a_1x+a_2x^2+\cdots
\]
に $x^n$ を掛けて、$n=0,1,2,\ldots$ について足し合わせると、
\[
  \sum_{n=0}^{\infty}a_{n+2}x^n
  =\sum_{n=0}^{\infty}a_{n+1}x^n
  +\sum_{n=0}^{\infty}a_nx^n
\]
となる。

それぞれを $A(x)$ で書き直すと、
\[
  \frac{A(x)-a_0-a_1x}{x^2}
  =\frac{A(x)-a_0}{x}+A(x)
\]
である。
整理すれば、
\[
  A(x)=\frac{a_0+(a_1-a_0)x}{1-x-x^2}
\]
を得る。

ここで重要なのは、漸化式の各項を追いかける代わりに、数列全体についての方程式を作れたことである。

\subsection{母関数から一般項を取り出す}

\[
  \varphi=\frac{1+\sqrt5}{2},\qquad
  \psi=\frac{1-\sqrt5}{2}
\]
とおくと、
\[
  1-x-x^2=(1-\varphi x)(1-\psi x)
\]
である。
また、形式的冪級数として、
\[
  \frac1{1-rx}=1+rx+r^2x^2+\cdots
\]
が成り立つ。

したがって、部分分数分解を使うと、
\[
  a_n=
  \frac{a_1-a_0\psi}{\sqrt5}\varphi^n
  +\frac{a_0\varphi-a_1}{\sqrt5}\psi^n
\]
となる。

母関数による解法の流れは、
\[
  \text{漸化式}
  \longrightarrow
  \text{母関数}
  \longrightarrow
  \text{分数式}
  \longrightarrow
  \text{係数}
\]
である。
本編で扱った変数変換とは違う道具だが、「扱いにくい対象を、扱いやすい対象へ移す」という発想は同じである。

\section{複数の数列を一つにまとめる}

\subsection{連立漸化式とベクトル}

二つの数列が互いに影響し合うと、
\[
  a_{n+1}=pa_n+qb_n,
  \qquad
  b_{n+1}=ra_n+sb_n
\]
のような連立漸化式になる。
二本の式を別々に扱う代わりに、縦に並べて一つのベクトルとみなす。

\[
  \boldsymbol v_n=
  \begin{pmatrix}a_n\\b_n\end{pmatrix},\qquad
  A=\begin{pmatrix}p&q\\r&s\end{pmatrix}
\]
とおけば、
\[
  \boldsymbol v_{n+1}=A\boldsymbol v_n
\]
となる。

これは、二つの数列を一つの更新規則で動かしているという意味である。
繰り返し代入すれば、
\[
  \boldsymbol v_n=A^n\boldsymbol v_0
\]
となる。
漸化式を解く問題が、行列の累乗を調べる問題へ変わった。

\subsection{フィボナッチ数列を行列で見る}

フィボナッチ数列を
\[
  F_0=0,\qquad F_1=1,\qquad
  F_{n+2}=F_{n+1}+F_n
\]
で定める。
二つの隣り合う項を並べると、
\[
  \begin{pmatrix}F_{n+1}\\F_n\end{pmatrix}
  =
  \begin{pmatrix}1&1\\1&0\end{pmatrix}^{n}
  \begin{pmatrix}1\\0\end{pmatrix}
\]
となる。

一見すると、行列は新しい道具を持ち込んだだけに見える。
しかし実際には、二つの項を一つの状態として記録しただけである。
「次の項を決めるには、直前の二項が必要」という情報を、そのままベクトルにしたのである。

\subsection{固有値が見つかると何がうれしいか}

行列 $A$ に対して、零ベクトルでない $\boldsymbol v$ が
\[
  A\boldsymbol v=\lambda\boldsymbol v
\]
を満たすとき、$\lambda$ を\textbf{固有値}、$\boldsymbol v$ を\textbf{固有ベクトル}という。

固有ベクトルの方向のベクトルは、$A$ を掛けても方向が変わらず、長さだけ $\lambda$ 倍になる。
したがって、その方向に分解できれば、行列の累乗は単なる数の累乗に変わる。

フィボナッチ行列
\[
  A=\begin{pmatrix}1&1\\1&0\end{pmatrix}
\]
の固有値は、
\[
  \lambda^2-\lambda-1=0
\]
の解、すなわち $\varphi,\psi$ である。
母関数から現れた数と、行列から現れた数が一致した。
これは偶然ではない。
どちらも、同じ二項間漸化式の構造を取り出しているからである。

\subsection{高階漸化式も行列にできる}

\[
  a_{n+2}=pa_{n+1}+qa_n
\]
に対して、
\[
  \begin{pmatrix}a_{n+2}\\a_{n+1}\end{pmatrix}
  =
  \begin{pmatrix}p&q\\1&0\end{pmatrix}
  \begin{pmatrix}a_{n+1}\\a_n\end{pmatrix}
\]
と書ける。
三項以上を使う漸化式でも、必要な項をベクトルにまとめれば同じ発想が使える。
特性方程式による解法と行列による解法は、違う計算をしているようで、実は同じ構造を別の言葉で語っているのである。

\section{三角関数と多項式をつなぐ}

\subsection{二項間漸化式から出発する}

次の漸化式を考える。
\[
  a_{n+1}=2x a_n-a_{n-1}
\]
ここで $x$ は固定された数である。
これは、本編で学んだ二項間漸化式の一種であり、特性方程式は、
\[
  r^2-2xr+1=0
\]
となる。

もし $-1\leq x\leq1$ ならば、$x=\cos\theta$ と書ける。
すると特性方程式の解は、複素数を使えば $e^{i\theta},e^{-i\theta}$ となる。
しかし、ここでは複素数を使わずに、三角関数の加法定理から同じ構造を見つけてみよう。

\[
  \cos((n+1)\theta)
  =2\cos\theta\cos(n\theta)-\cos((n-1)\theta)
\]
だから、$\cos(n\theta)$ は、
\[
  u_{n+1}=2\cos\theta\,u_n-u_{n-1}
\]
を満たしている。

つまり、三角関数の倍角・加法の世界と、二項間漸化式の世界が同じ規則を持っている。
この対応を多項式として記録したものが、チェビシェフ多項式である。

\subsection{チェビシェフ多項式の定義}

第一種チェビシェフ多項式 $T_n(x)$ を、
\[
  T_0(x)=1,\qquad T_1(x)=x,\qquad
  T_{n+1}(x)=2xT_n(x)-T_{n-1}(x)
\]
で定める。

この定義から、
\[
  T_2(x)=2x^2-1,\qquad
  T_3(x)=4x^3-3x
\]
を得る。
漸化式の右辺が多項式だけでできているので、$T_n(x)$ は確かに多項式である。

一方、$x=\cos\theta$ と置くと、三角関数の関係から、
\[
  T_n(\cos\theta)=\cos(n\theta)
\]
が成り立つ。
実際、$n=0,1$ では定義から明らかであり、$n$ と $n-1$ で成り立つと仮定すれば、
\[
\begin{aligned}
  T_{n+1}(\cos\theta)
  &=2\cos\theta\,T_n(\cos\theta)-T_{n-1}(\cos\theta)\\
  &=2\cos\theta\cos(n\theta)-\cos((n-1)\theta)\\
  &=\cos((n+1)\theta)
\end{aligned}
\]
となる。

これは、三角関数を多項式の言葉で表す公式である。
\[
  \cos(2\theta)=2\cos^2\theta-1
\]
という倍角公式も、$T_2(x)=2x^2-1$ と書けば、同じ一つの仕組みとして見えてくる。

\subsection{この話が漸化式の本にある理由}

チェビシェフ多項式は、名前を覚えるために登場したのではない。
同じ漸化式から、次の四つの見方が現れることが重要である。

\begin{itemize}
  \item 特性方程式で解く。
  \item $x=\cos\theta$ と置いて三角関数で表す。
  \item 行列の固有値として理解する。
  \item 多項式 $T_n(x)$ として、変数 $x$ の関数にする。
\end{itemize}

一つの漸化式が、代数・三角関数・行列・多項式を結ぶ橋になっている。
本編で学んだ「式の形を見る」という態度が、ここでは数学の分野同士をつなぐ視点になっている。

\section{一般項の先を見る：漸近解析}

\subsection{一般項がなくても知りたいこと}

漸化式を解くとき、一般項を求めることは重要である。
しかし、一般項が初等的な式で書けない場合でも、数列について知りたいことは残っている。

例えば、
\begin{itemize}
  \item 数列は増えているのか、減っているのか。
  \item 上から押さえられているのか。
  \item $n$ が大きくなると、どのくらいの速さで増えるのか。
  \item 極限を持つのか。
\end{itemize}

このように、大きな $n$ における振る舞いを調べることを\textbf{漸近解析}という。

\subsection{極限があるなら不動点になる}

\[
  a_{n+1}=f(a_n)
\]
に対して、もし $a_n\to L$ ならば、通常は $a_{n+1}\to L$ でもある。
したがって、関数が連続なら、
\[
  L=f(L)
\]
が必要になる。
この式を解けば、極限の候補を絞ることができる。

ただし、前に見た
\[
  a_{n+1}=2-a_n
\]
のように、不動点があっても収束しない場合がある。
不動点の周りで値が縮むのか、跳ね返るのかを調べなければならない。
関数 $f$ が微分可能で、$|f'(L)|<1$ ならば、$L$ の近くの値は縮みながら近づきやすい。
逆に $|f'(L)|>1$ ならば、少しのずれが広がりやすい。
これは厳密な定理を学ぶ前の、動きを読むための目安である。

\subsection{増加の速さを比べる}

数列の値を一つずつ求めなくても、増加の速さを比較できる。
例えば、十分大きな $n$ では、
\[
  n<n^2<2^n
\]
となる。
多項式的な増加よりも、指数的な増加の方が速いのである。

二つの数列 $a_n,b_n$ について、
\[
  \lim_{n\to\infty}\frac{a_n}{b_n}=1
\]
なら、
\[
  a_n\sim b_n
\]
と書く。
これは、二つの数列の差が小さいというより、比が $1$ に近づくという意味である。

フィボナッチ数列では、
\[
  F_n\sim\frac1{\sqrt5}
  \left(\frac{1+\sqrt5}{2}\right)^n
\]
となる。
正確な一般項を知るだけでなく、「どの項が増加の中心を担っているか」を知ることができる。

\section{「解けない」から始まる数学}

\subsection{この本でいう「解けない」}

本書では、漸化式を「解ける」とは、主に高校数学で扱う関数や有限回の操作によって一般項を表せることだと考えてきた。
したがって、「解けない」とは、数学的に何も分からないという意味ではない。
本書の方法では、初等的な一般項を作ることが難しいという意味である。

例えば、
\[
  a_{n+1}=a_n^2+1,\qquad a_0=0
\]
なら、
\[
  0,1,2,5,26,677,\ldots
\]
と計算できる。
一般項を簡単な式で書けなくても、項を計算すること、増加の速さを調べること、正負や単調性を調べることはできる。

ここで大切なのは、一般項を求めることだけを数学の終点にしないことである。
問題によっては、
\[
  \text{正確な値}
\]
よりも、
\[
  \text{動き方、増え方、安定性、極限}
\]
を知る方が本質的なこともある。

\subsection{漸化式から広がった地図}

この付録で見てきたことを、一つの流れに戻してみよう。

まず、隣り合う項の差を調べた。
そこから、連続的な変化を記述する微分と微分方程式へ進んだ。
次に、解を直接求められないとき、解に近づく数列を作るという発想を見た。
さらに、数列全体を一つの式にまとめる母関数、複数の項を一つのベクトルにまとめる行列へ進んだ。

二項間漸化式を三角関数で眺めると、チェビシェフ多項式が現れた。
そこでは、一つの漸化式が、多項式・三角関数・行列・特性方程式をつないでいた。
最後に、一般項そのものではなく、数列の大きな振る舞いを調べる漸近解析へたどり着いた。

\[
\begin{array}{c}
\text{隣り合う項の関係}\\
\downarrow\\
\text{差分・微分}\\
\downarrow\\
\text{連続化・数値計算}\\
\downarrow\\
\text{反復・近似}\\
\downarrow\\
\text{母関数・行列}\\
\downarrow\\
\text{三角関数・多項式}\\
\downarrow\\
\text{極限・漸近解析}
\end{array}
\]

漸化式は、最初は数列を計算するための短い式に見える。
しかし、その中には「変化を記録する」「同じ操作を繰り返す」「複雑なものを別の形に移す」「大きな動きを予測する」という、数学の多くの考え方が含まれている。

漸化式を解くとは、一般項を見つけることだけではない。
その式が何を記録し、どのような構造を持ち、どの数学へつながっているのかを見つけることでもある。
